% Generated draft: requires mathematical and source review.
\documentclass[11pt]{article}
\usepackage[margin=25mm]{geometry}
\usepackage{fontspec}
\setmainfont{Latin Modern Roman}
\usepackage{amsmath,amssymb,mathtools,unicode-math}
\setmathfont{Latin Modern Math}
\usepackage{enumitem,needspace,xurl,hyperref}
\setlength{\parindent}{0pt}
\setlength{\parskip}{5pt}
\setlength{\emergencystretch}{3em}
\newcommand{\N}{\mathbb{N}}
\newcommand{\Z}{\mathbb{Z}}
\newcommand{\Q}{\mathbb{Q}}
\newcommand{\R}{\mathbb{R}}
\newcommand{\C}{\mathbb{C}}
\newcommand{\F}{\mathbb{F}}
\newcommand{\A}{\mathbb{A}}
\newcommand{\PP}{\mathbb{P}}
\newcommand{\E}{\mathbb{E}}
\newcommand{\eps}{\varepsilon}
\newcommand{\dd}{\,\mathrm{d}}
\begin{document}
\section*{UnsolvedMath ID: 801023}
\textbf{Source problem code:} catalogue entry
\subsection{Definitions and mathematical statement}
A profinite group is a topological group isomorphic to an inverse limit of finite discrete groups. It is topologically finitely generated if there are finitely many elements \(g_1,\ldots,g_m\in G\) such that
\[
G=\overline{\langle g_1,\ldots,g_m\rangle}.
\]
Under the usual lower-cut convention, a real \(r\) is left \(\Sigma^0_2\) if its strict rational lower cut \(\{q\in\Q:q<r\}\), under an effective coding of \(\Q\), is a \(\Sigma^0_2\) set. The supplied record uses the term ``left \(\Sigma_2\)'' without defining it.

For a specified metric \(d\), write \(\dim_H^d(H)\) for the Hausdorff dimension of \(H\) with the restricted metric. The source states, as background, that every left \(\Sigma_2\) real \(r\in[0,1]\) is realized as the Hausdorff dimension of a computable subgroup of a computable profinite group.

The question is to determine the attainable values of \(r\) when the ambient group is additionally topologically finitely generated. Conditional on fixing the source-intended effective presentations and its metric assignment \(G\mapsto d_G\), the requested set can be written as
\[
\mathcal R_{\mathrm{fg}}=
\left\{r\in[0,1]:
\begin{gathered}
\text{there exist a computable profinite group }G\text{ and a computable subgroup }H\leq G,\\
G\text{ is topologically finitely generated},\qquad
\dim_H^{d_G}(H)=r
\end{gathered}
\right\}.
\]
This formulation remains conditional: the record does not specify the meaning of computability for either object, the metric or filtration defining Hausdorff dimension, or any normalization of that dimension. It also does not specify whether the subgroup must be closed. These choices are required to make the attainable-value question precise. The wording does not explicitly resolve whether the requested values are restricted in advance to left \(\Sigma_2\) reals or range over all of \([0,1]\).
\subsection{Short English statement}
The source asks which Hausdorff dimensions can occur for computable subgroups of computable profinite groups when the ambient group has finitely many topological generators. It supplies a realization assertion for left \(\Sigma_2\) reals in \([0,1]\) without the finite-generation requirement. The effective presentation, dimension metric, subgroup closedness, and intended range of candidate reals require clarification.
\subsection{Sources}
\begin{itemize}
\item Local UnsolvedMath registry: supplied complete source record, permanent ID \(801023\), exact source code \texttt{catalogue entry}, entitled \emph{Open Problems in Mathematical Logic --- Problem 35}.
\item Supplied attribution: Barmpalias, Gao et al., \emph{Open Problems in Mathematical Logic} (2026), Problem 35. Supplied source URL: \url{https://arxiv.org/abs/2608.26628}. The record reports access on 2026-10-16 and extraction by \texttt{arxiv\_latex}; these are supplied provenance details, not independently verified access.
\item The registry labels this problem ``open.'' This is a supplied status label, not an independently verified mathematical or literature assessment. The background realization assertion is likewise reproduced as a source assertion.
\end{itemize}
\end{document}
